% Options for packages loaded elsewhere
\PassOptionsToPackage{unicode}{hyperref}
\PassOptionsToPackage{hyphens}{url}
\documentclass[
  ignorenonframetext,
]{beamer}
\newif\ifbibliography
\usepackage{pgfpages}
\setbeamertemplate{caption}[numbered]
\setbeamertemplate{caption label separator}{: }
\setbeamercolor{caption name}{fg=normal text.fg}
\beamertemplatenavigationsymbolsempty
% remove section numbering
\setbeamertemplate{part page}{
  \centering
  \begin{beamercolorbox}[sep=16pt,center]{part title}
    \usebeamerfont{part title}\insertpart\par
  \end{beamercolorbox}
}
\setbeamertemplate{section page}{
  \centering
  \begin{beamercolorbox}[sep=12pt,center]{section title}
    \usebeamerfont{section title}\insertsection\par
  \end{beamercolorbox}
}
\setbeamertemplate{subsection page}{
  \centering
  \begin{beamercolorbox}[sep=8pt,center]{subsection title}
    \usebeamerfont{subsection title}\insertsubsection\par
  \end{beamercolorbox}
}
% Prevent slide breaks in the middle of a paragraph
\widowpenalties 1 10000
\raggedbottom
\AtBeginPart{
  \frame{\partpage}
}
\AtBeginSection{
  \ifbibliography
  \else
    \frame{\sectionpage}
  \fi
}
\AtBeginSubsection{
  \frame{\subsectionpage}
}
\usepackage{iftex}
\ifPDFTeX
  \usepackage[T1]{fontenc}
  \usepackage[utf8]{inputenc}
  \usepackage{textcomp} % provide euro and other symbols
\else % if luatex or xetex
  \usepackage{unicode-math} % this also loads fontspec
  \defaultfontfeatures{Scale=MatchLowercase}
  \defaultfontfeatures[\rmfamily]{Ligatures=TeX,Scale=1}
\fi
\usepackage{lmodern}
\ifPDFTeX\else
  % xetex/luatex font selection
\fi
% Use upquote if available, for straight quotes in verbatim environments
\IfFileExists{upquote.sty}{\usepackage{upquote}}{}
\IfFileExists{microtype.sty}{% use microtype if available
  \usepackage[]{microtype}
  \UseMicrotypeSet[protrusion]{basicmath} % disable protrusion for tt fonts
}{}
\makeatletter
\@ifundefined{KOMAClassName}{% if non-KOMA class
  \IfFileExists{parskip.sty}{%
    \usepackage{parskip}
  }{% else
    \setlength{\parindent}{0pt}
    \setlength{\parskip}{6pt plus 2pt minus 1pt}}
}{% if KOMA class
  \KOMAoptions{parskip=half}}
\makeatother
\setlength{\emergencystretch}{3em} % prevent overfull lines
\providecommand{\tightlist}{%
  \setlength{\itemsep}{0pt}\setlength{\parskip}{0pt}}
\usepackage{bookmark}
\IfFileExists{xurl.sty}{\usepackage{xurl}}{} % add URL line breaks if available
\urlstyle{same}
\hypersetup{
  hidelinks,
  pdfcreator={LaTeX via pandoc}}

\author{\texorpdfstring{}{}}
\date{}

\begin{document}

\begin{frame}[fragile]{Introduction}
\protect\phantomsection\label{introduction}
This \texttt{code\_project} serves as the foundational exemplar for the
\href{https://github.com/docxology/template}{docxology/template}
ecosystem, demonstrating a fully-tested numerical optimization
implementation securely bracketed by rigorous infrastructure, hermetic
testing, and extensive documentation architectures. The words you are
reading---and the mathematical figures below them---have been
programmably generated through an unbreakable custody chain starting
from algorithm implementation through strict CI/CD validation to
multi-format \texttt{.pdf} compilation.

\begin{block}{Template Architecture Context}
\protect\phantomsection\label{template-architecture-context}
Scientific engineering requires mathematical accuracy combined with
software reliability. This project unifies theoretical optimization with
the repository's three foundational pillars:

\begin{enumerate}
\tightlist
\item
  \textbf{\texttt{infrastructure/} Layer (Root Directory)}: A modular
  stack of nine subpackages (\texttt{core}, \texttt{validation},
  \texttt{scientific}, \texttt{rendering}, \texttt{reporting},
  \texttt{documentation}, \texttt{llm}, \texttt{publishing},
  \texttt{project}) providing the computational scaffolding.
\item
  \textbf{\texttt{tests/} Framework
  (\texttt{projects/code\_project/tests/})}: An uncompromising
  validation layer maintaining a zero-mock testing policy. This is
  enforced automatically via
  \href{https://github.com/docxology/template/blob/main/.github/workflows/ci.yml}{\texttt{.github/workflows/ci.yml}}
  mapping to \texttt{pyproject.toml} directives.
\item
  \textbf{\texttt{docs/} Knowledge Base
  (\texttt{projects/code\_project/docs/})}: A structured repository of
  architectural guidelines, operational patterns, and the Rigorous
  Agentic Scientific Protocol (RASP) that governs the AI-assisted agents
  writing these very texts.
\end{enumerate}

This implementation of gradient descent algorithms for solving
optimization problems is used as the vehicle to demonstrate these
pillars. The theoretical problem is mapped programmatically inside
\href{https://github.com/docxology/template/blob/main/projects/code_project/src/optimizer.py}{\texttt{projects/code\_project/src/optimizer.py}}:

\begin{equation}
\label{eq:optimization_problem}
\min_{x \in \mathbb{R}^n} f(x)
\end{equation}

where \(f: \mathbb{R}^n \rightarrow \mathbb{R}\) is a continuously
differentiable objective function.
\end{block}

\begin{block}{Infrastructure Integration}
\protect\phantomsection\label{infrastructure-integration}
Rather than existing as isolated scripts, this project extensively
leverages the \texttt{infrastructure} layer:

\begin{itemize}
\tightlist
\item
  \textbf{Scientific Utilities}: Utilizing
  \texttt{infrastructure.scientific.stability} and
  \texttt{infrastructure.scientific.benchmarking} to guarantee numerical
  boundaries and performance scaling.
\item
  \textbf{Hermetic Validation}: Deploying
  \texttt{infrastructure.validation} components
  (\texttt{markdown\_validator}, \texttt{output\_validator}) to ensure
  all generated artifacts are cryptographically sound.
\item
  \textbf{Reporting \& Rendering}: Employing
  \texttt{infrastructure.rendering.pdf\_renderer} and
  \texttt{infrastructure.reporting.executive\_reporter} to automatically
  transform code outputs into this finalized manuscript.
\end{itemize}
\end{block}

\begin{block}{Algorithm Overview}
\protect\phantomsection\label{algorithm-overview}
The reference gradient descent algorithm iteratively updates the
solution using:

\begin{equation}
\label{eq:gradient_descent_update}
x_{k+1} = x_k - \alpha \nabla f(x_k)
\end{equation}

where \(\alpha > 0\) is the step size (learning rate) and
\(\nabla f(x_k)\) is the gradient of the objective function at iteration
\(k\).
\end{block}

\begin{block}{Exemplar Implementation Goals}
\protect\phantomsection\label{exemplar-implementation-goals}
As the representative project for the repository, this implementation
explicitly demonstrates:

\begin{enumerate}
\tightlist
\item
  \textbf{Infrastructure-Coupled Code}: Scientific implementations that
  delegate logging, file ops, and reporting to the
  \texttt{infrastructure} core.
\item
  \textbf{Zero-Mock Verification}: A strict 34-test validation suite
  proving numerical accuracy without artificial test boundaries.
\item
  \textbf{Automated Research Pipelines}: High-precision analyses that
  generate publication-quality, accessible visualizations automatically.
\item
  \textbf{Agentic Documentation standards}: Native adherence to the RASP
  methodology and \texttt{AGENTS.md} guidelines, ensuring the logic
  remains verifiable by both human and artificial intelligence.
\end{enumerate}
\end{block}
\end{frame}

\end{document}
